\documentclass{article}
\usepackage[utf8]{inputenc}
\usepackage[T1]{fontenc}
\usepackage[a4paper,total={6in,8in}]{geometry}
\usepackage{amsmath,amssymb,amsfonts}
\usepackage{graphicx}
\usepackage{hyperref}
\usepackage{enumitem}
\usepackage{listings}
\usepackage{xcolor}
\DeclareMathOperator{\Var}{Var}
\DeclareMathOperator{\Cov}{Cov}
\providecommand{\tightlist}{\setlength{\itemsep}{0pt}\setlength{\parskip}{0pt}}
\lstset{basicstyle=\ttfamily\scriptsize,breaklines=true,columns=fullflexible,
  keepspaces=true,showstringspaces=false,frame=single}
\hypersetup{colorlinks=true,urlcolor=blue,linkcolor=blue}
\setlength{\emergencystretch}{2em}

\title{Math 126 Problem 19: Least Squares Regression}
\author{} % Put your name here.
\date{Spring 2027}
\begin{document}
\maketitle

\hypertarget{submission-instructions}{%
\subsection*{Submission Instructions}\label{submission-instructions}}

Please upload your work to Gradescope by \textbf{11:59pm on Monday, April 19, 2027 (Pacific time).}

\begin{itemize}
\item
  You are encouraged to discuss ideas and work with your classmates. However, you \emph{must acknowledge} your collaborators at the top of each solution on which you collaborated with others and you \emph{must write} your solutions independently.
\item
  Your solutions to theory questions must be written legibly, or typeset in LaTeX or markdown.
\item
  I recommend that you write your solutions in a Python notebook using Google Colab.
\item
  You should submit your solutions as a \emph{single PDF} via the assignment on Gradescope.
\end{itemize}

\textbf{Grading:} The point of the problem set is for \emph{you} to learn. To this end, I hope to disincentivize the use of LLMs by \emph{not} grading your work for correctness. Instead, you will grade your own work by comparing it to my solutions. This self-grade is due the Friday \emph{after} the problem set is due, also on Gradescope.

\noindent Original notebook: \href{https://colab.research.google.com/drive/19zPXyyDEFFtQ04rdEqAKh5Osd2ybwI1S?usp=sharing}{Google Colab}.

% Source notebook cell 2
\hypertarget{least-squares-regression}{%
\subsection*{Least Squares Regression}\label{least-squares-regression}}

Let \(n \geq d\). Consider a full-rank, feature matrix \(\mathbf{A} \in \mathbb{R}^{n \times d}\) and a target vector \(\mathbf{b} \in \mathbb{R}^n\). The regression problem is to find a minimizing vector \begin{align*}
\mathbf{x}^* = \arg \min_{\mathbf{x} \in \mathbb{R}^d}
\| \mathbf{Ax} - \mathbf{b} \|_2^2.
\end{align*}

In class, we showed how to quickly solve the regression problem by projecting our feature matrix and target vector into a lower dimensional space. In this problem, you will derive the exact solution for \(\mathbf{x}^*\) and explore how long it takes to compute it.

\hypertarget{part-a-optimal-solution}{%
\subsection*{Part A: Optimal Solution}\label{part-a-optimal-solution}}

Show that \(\mathbf{x}^* = \left(\mathbf{A}^\top \mathbf{A} \right)^{-1} \mathbf{A}^\top \mathbf{b}\).

\textbf{Hint:} Compute the partial derivative with respect to \(x_k\) of the \(\ell_2\)-norm and set it equal to 0.

\hypertarget{part-b-time-complexity}{%
\subsection*{Part B: Time Complexity}\label{part-b-time-complexity}}

Explain why computing \(\left(\mathbf{A}^\top \mathbf{A} \right)^{-1} \mathbf{A}^\top \mathbf{b}\) takes \(O(nd^2)\) time. Then, load a regression dataset of your choice (e.g., from \texttt{scikit-learn}), build \(\mathbf{A}\) and \(\mathbf{b}\) and find the optimal solution \(\mathbf{x}^*\).

\end{document}
